\section{Mathematical Foundations}

\begin{purposebox}
	\textbf{Purpose:} Demonstrate understanding of the core mathematics needed to understand the theoretical framework of the tracking control NN.
	
	\textbf{When complete:} You can confidently read the paper's equations and trace each step of the derivations.
\end{purposebox}

%──────────────────────────────────────────────────
\subsection{Difference Equations}
%──────────────────────────────────────────────────

\subsubsection*{Concept Check}

\question{Basic Form \& Solutions}
Write the general form of a first-order difference equation. Then, given $x(k+1) = 2x(k) + 1$ with $x(0) = 0$, solve for $x(1),\; x(2),\; x(3),\; x(4)$ by hand.
\begin{answerbox}
	\centering General Form: $x_{ k+1 } = F(x_{ k })$
\end{answerbox}
\begin{answerbox}
	\centering Given: $x(k+1) = 2x(k) + 1$, $x(0) = 0$
	\begin{align}
		x(0) = 0                 \\
		x(1) & = 2x(0) + 1 = 1,  \\
		x(2) & = 2x(1) + 1 = 3   \\
		x(3) & = 2x(2) + 1 = 7   \\
		x(4) & = 2x(3) + 1 = 15
	\end{align}
\end{answerbox}

\question{Variable Structure Systems}
In your own words, explain what makes a system ``variable structure.'' Why does the signum function $\operatorname{sgn}(e)$ create variable structure behavior in the tracking control NN?

\begin{answerbox}
	A system is of variable structure if the closed loop update law changes discontinuously
	based on certain conditions of the input i.e. modified by a signum function.
\end{answerbox}

\question{Euler Approximation (Discretization)}
Given the continuous-time ODE $\dot{x} = f(x)$, derive the Euler approximation that produces the discrete-time form $x(k+1) = x(k) + \tau f(x(k))$. What role does $\tau$ play, and what happens if $\tau$ is too large?

\begin{answerbox}
	\begin{math}
		\begin{aligned}
			\text{given $x\prime = f(x)$:}                                                                           \\
			x\prime                                             & = f(x)                                             \\
			x\prime(t) \approx  \frac{x(t + \tau) - x(k)}{\tau} & \approx f(x(t))                                    \\
			x(t + \tau) - x(t)                                  & \approx \tau( f(x(t)) )                            \\
			x(t + \tau)                                         & \approx x(t) + \tau( f(x(t)) )                     \\
			x(k + 1)                                            & \approx x(k) + \tau( f(x(k)) )\text{, $k = t\tau$} \\
		\end{aligned}
	\end{math} \\
	Tau is the quotient of the current time step over
	the time horizon(total amount of time steps)
\end{answerbox}

\question{Connecting to the Paper}
The paper's core equation is:
\[
	x(k+1) = x(k) + \tau\bigl[f(x(k)) + g(x(k))\,u(k)\bigr]
\]
Identify which part comes from the Euler approximation of a
continuous-time system, and write out what the original
continuous-time ODE would look like.

\begin{answerbox}
	$x(k+1), x(k), and \tau$ come from the Euler approximation. \\
	The original ODE:
	\begin{equation}
		\dot{x}(t) = f(x(t)) + g(x(t))\,u(t)
	\end{equation} \\ \\
\end{answerbox}

%──────────────────────────────────────────────────
\subsection{Linear Algebra}
%──────────────────────────────────────────────────

\subsubsection*{Concept Check}

\question{Matrix Operations}
Given
\[
	A = \begin{bmatrix} 2 & 1 \\ 0 & 3 \end{bmatrix}, \qquad
	B = \begin{bmatrix} 1 & 0 \\ 2 & 1 \end{bmatrix},
\]
compute $AB$, $A^\top$, and $A^{-1}$. Show your work.

\begin{largeanswer}\end{largeanswer}

\question{Eigenvalues \& Stability}
Compute the eigenvalues of $A = \begin{bmatrix} 2 & 1 \\ 0 & 3 \end{bmatrix}$. For a discrete-time linear system $x(k+1) = Ax(k)$, what condition on the eigenvalues guarantees stability? Is this system stable?

\begin{answerbox}\end{answerbox}

\question{Diagonal Matrices}
The paper uses $\beta = \operatorname{diag}(\beta_1, \dots, \beta_n)$ as the learning rate parameter matrix. Explain what a diagonal matrix is, and why using a diagonal matrix for $\beta$ means each state variable has its own independent learning rate.

\begin{answerbox}\end{answerbox}

%──────────────────────────────────────────────────
\subsection{Optimization \& Gradient}
%──────────────────────────────────────────────────

\subsubsection*{Concept Check}

\question{Gradient Descent}
Explain the gradient descent update rule $\Delta e = -\nabla C(e)$ in plain language. Why does moving in the negative gradient direction minimize the cost?

\begin{answerbox}\end{answerbox}

\question{Cost Function Design}
The paper uses $C(x) = |e|$ where $e$ is the tracking error. Why is absolute value a valid cost function? What property does it have at $e = 0$? How does this connect to the signum function (hint: $\tfrac{d|e|}{de} = \operatorname{sgn}(e)$)?

\begin{answerbox}\end{answerbox}

\question{Convexity}
Define what it means for a cost function to be convex. Why does convexity guarantee a unique global minimum? Is $C(x) = |e|$ convex?

\begin{answerbox}\end{answerbox}

%──────────────────────────────────────────────────
\subsection*{Self-Assessment Checklist}
%──────────────────────────────────────────────────

\begin{itemize}[leftmargin=2em, label=$\square$]
	\item I can solve difference equations iteratively given initial conditions
	\item I can explain why the signum function creates variable structure
	\item I can derive the Euler approximation from a continuous ODE
	\item I can perform basic matrix operations (multiply, transpose, inverse)
	\item I can compute eigenvalues and relate them to system stability
	\item I can explain gradient descent and why the cost function $C = |e|$ is chosen
	\item I can connect all of the above to the paper's system model equation
\end{itemize}

\begin{notesbox}\end{notesbox}
